% Created 2026-09-26 Sat 23:39
% Intended LaTeX compiler: xelatex
\documentclass{article}

\usepackage{amsmath}
\usepackage{fontspec}
\usepackage{graphicx}
\usepackage{longtable}
\usepackage{wrapfig}
\usepackage{rotating}
\usepackage[normalem]{ulem}
\usepackage{capt-of}
\usepackage{hyperref}
\newcommand{\avg}[1]{\langle #1 \rangle}
\usepackage{amsmath}
\usepackage{amssymb}
\usepackage[thinc]{physics}
\usepackage{geometry}
\geometry{left=2.5cm,right=2.5cm,top=2.5cm,bottom=2.5cm}
\usepackage{xcolor}
\usepackage[unicode]{hyperref}
\hypersetup{
colorlinks=true,
linkcolor=blue,
citecolor=blue,
urlcolor=blue
}
\urlstyle{same} % keep URL font same as main text
\author{Soroush Seirani G.}
\date{\today}
\title{Avalanche's Progress Report 6}
\hypersetup{
 pdfauthor={Soroush Seirani G.},
 pdftitle={Avalanche's Progress Report 6},
 pdfkeywords={},
 pdfsubject={},
 pdfcreator={},
 pdflang={English}}
\begin{document}

\maketitle
\section{Optimization}
\label{sec:org645a572}
I wrote an optimization backend to find for a variable parameter \(N\),
which \(D\), \(\tau\) and \(R\) give us the least-square error for two given
simulation results (probability distributions).

\begin{align}
  x[i] &\longrightarrow \frac{x[i]}{N^D}\\
  y[i] &\longrightarrow y[i] \cdot x^{\tau}[i]
\end{align}

However due to the sheer amount of compute per iteration this wasn't
feasible and didn't produce good results. I did manage to write a
single-parameter optimization which only optimized for \(R\). However, I
didn't get to high values of \(R\) because supercritical systems have a
long simulation time. This made me look into a way to speed up the
simulation.
\section{Algorithmic Optimization of the Stochastic Simulation}
\label{sec:org58d5535}
Since I had used up all the numerical tricks I knew of, I turned to an
LLM to get some ideas. It told me it's common practice to model this
avalanche as a branching process. It offered an algorithm of time
complexity \(O(N)\). I spent most of my time studying this new
algorithm. I've adapted an elementary version of this which isn't
complete (isn't of \(O(N)\) yet) but I'll complete it in the following
week.

The basic idea is to think of avalanches as the transitions that
happen between levels. Next we should realize that the order of
transitions doesn't matter. By interchanging the order of transitions,
neither the probability nor the size change. So we can reorder the
transitions into a desired format which is easy for us to analyze.

Let \(A\) be the number of active neurons and also designates the level
we're at. Being at level \(A\) means there are \(A\) active neurons. Let
\(D_K\) be the number of \emph{downcrossings}: The number of times we go from
\(k\) to \(k-1\). By this definition \(D_1 = 1\) all the time.

\begin{gather}
  D_k := k \rightarrow k-1
\end{gather}

We have also defined \(p_k\) and \(q_k\):

\begin{gather}
  p_k := \text{probability of } k \rightarrow k-1\\
  q_k := \text{probability of } k \rightarrow k+1
\end{gather}

Now let's assume we're at level \(k\): There is the probability of going
down and going up. We call the probability of going down the first
time:

\begin{gather}
  G_k(0) = p_k
\end{gather}

Where \(G_k(a)\) denotes the probability we will go from \(k\) to \(k+1\) before
going down to \(k-1\):

\begin{gather}
  G_k(1) = q_kp_k\\
  G_k(a) = q_k^ap_k
\end{gather}

What I'm trying to get across is this: \textbf{\(G_k(a)\) denotes the
probability of going up from level \(k\) \(a\) times before finally going
down to \(k-1\)}.

It's important to see it's normalized:

\begin{align}
\sum_{i=0}^\infty G_k(i) &= p_k\sum_{i=0}^\infty q^i_k \text{ since } q_k<1 \\
  &= p_k \frac{1}{1-q_k} = 1
\end{align}

Now let's look at level 2. Since we know \(D_1=1\) we can see:

\begin{gather}
  D_2 = G_1
\end{gather}

Where \(G_1\) denotes a random variable with the discrete probability
distribution \(P(G_1=a) = q_1^a p_1\). Now if \(D_2 > 2\) (say \(r\) for our
example) for \(D_3\) we'll have:

\begin{gather}
  D_3 = \sum_{i=1}^r G_{2,i}
\end{gather}

Where each of the \(G_{2,i}\) is a random variable with the probability
distribution of \(G_2\). We can generalize this to an arbitrary level
\(k\):

\begin{gather}
D_k = \sum_{i=1}^r G_{k-1,i}
\end{gather}

This is as far as I've managed to implement. This method gives the same
simulation results as the Gillespie algorithm. However we haven't yet
arrived at the time complexity of \(O(N)\). That's due to the arbitrary
sum of \(r\) random variables. At higher levels \(r\) can grow
increasingly large.

Since the sum of random geometric variables gives a negative binomial
distribution, by directly sampling such a distribution for \(D_k\) this
algorithm arrives exactly at \(O(N)\).

This is a fascinating result but since I'm not sure the approach is
correct I haven't applied it. My main concern is that it assumes all
avalanches die. On one hand we were already throwing out undying
avalanches but Once I'm certain the theory behind this algorithm is
correct I'll implement it properly.
\end{document}
